%!TEX root = ../main.tex
% =============================================================================
%  五、问题 2：不确定性下的两阶段计划购电模型      目标 3.5 页
%  评审在问：不确定性处理合理吗？
%  本节含：① 两阶段结构与非预期性（式 (12) 全式）
%          ② 场景集生成与削减（分层 + medoid，一次存盘固定）
%          ③ 报童临界分位数与两条界，及四条适用边界
%          ④ 对照：R1／R2／R3、无安全边际消融、机会约束、CVaR 扫描
%          ⑤ 结果与可行性复核 → 题面格式表见第八章
%  口径：C-05 紧急购电量在本节首次定义；图名 fig07_q2_emergency.pdf
% =============================================================================

\section{问题 2：不确定性下的两阶段计划购电模型}\label{chap:q2}

\subsection{模型建立}

\subsubsection{决策结构与信息结构}\label{sec:q2-struct}

与问题 1 的差别有三点：电价仍逐日相同，而负载与光伏\textbf{逐日逐时段变化}；计划仍在每天
\(0{:}00\) 制定；若微网供给低于负载，须按\textbf{交易时刻电价的 5 倍}紧急购电。故决策分两层：
第一阶段在 \(0{:}00\) 只依据预报一次性确定全天 144 槽的\textbf{计划购电量 \(E_t\)}与
\textbf{充放电量 \(c_t,d_t\)}，该决策\textbf{不依赖场景}；第二阶段在实际值揭示后，计划与实际之差由
\textbf{紧急购电量 \(g_{t,\omega}\)}（缺口，按 \(5p_t\) 计价）与盈余处置吸收。

信息结构取 \textbf{R3}：计划依据为附件 3 中当日 \(0{:}00\) 发布的光伏预报，题面未给负载预报，
故计划阶段按“当日负载曲线已知”处理（等价于负载预报无误差，假设 A9）；\textbf{结算依据为附件 2 的
当日实际负载与实际光伏}。两者是两套数据，其差异即紧急购电的来源。紧急购电不是决策量，而是由缺口
唯一确定的\textbf{结算量}。

\textbf{本节最关键、也最需要被检验的一条设定}（\textbf{假设 A12}）：充放电计划同样在第一阶段确定、
\textbf{实时不允许重调度}，故第一阶段买到的电量不能转手，超购部分既不产生收益也不再被后续时段使用
（\textbf{超购即弃}），其单位成本为 \(p_t\)。这一设定使模型天然具有\textbf{安全边际}，是报童型权衡
在本问成立的结构性原因；但它同时\emph{决定}了本文两项标志性输出的量级——紧急购电量
\(172\,015.82\)~kWh 与盈余处置量 \(3\,106\,390.84\)~kWh。题面只锁定\emph{结算}口径、未禁止储能按
实际值动作，故本文在 表~\ref{tab:abl-q2} 的 RT 行给出\textbf{实时再调度消融}：
固定已结算的 \(E^{\text{plan}}\)、用当日实际净需求重优化 \(c,d,S\)，量化该设定的代价。

\subsubsection{两阶段随机规划模型}\label{sec:q2-model}

\begin{subequations}\label{eq:q2}
\begin{align}
  \min_{E,c,d,S}\quad
  & \sum_{t}p_tE_t+\sum_{\omega\in\Omega}W_\omega\sum_{t}5p_t\,g_{t,\omega} \label{eq:q2-obj}\\
  \text{s.t.}\quad
  & E_t,\,c_t,\,d_t,\,S_t\ \text{与}\ \omega\ \text{无关}, && \forall t \label{eq:q2-nonanticip}\\
  & E_t-c_t+d_t\ \ge\ \bigl(L^{\text{fcst}}_t-\widehat{PV}^{\text{fcst}}_t\bigr)\Delta t, && \forall t \label{eq:q2-plan}\\
  & g_{t,\omega}\ \ge\ \bigl(L_{\omega,t}-PV_{\omega,t}\bigr)\Delta t+c_t-d_t-E_t,\quad g_{t,\omega}\ge0, && \forall t,\omega \label{eq:q2-gap}\\
  & S_t=S_{t-1}+\eta_cc_t-d_t/\eta_d,\quad \underline S\le S_t\le\overline S,\quad 0\le c_t,d_t\le\overline P\Delta t, && \forall t \label{eq:q2-soc}\\
  & E_t\ge0,\quad S_0=S^{init},\quad S_{144}\ge S_0 . && \label{eq:q2-nonneg}
\end{align}
\end{subequations}

式~\eqref{eq:q2-nonanticip} 是两阶段结构的代数表达：带 \(\omega\) 下标者为\textbf{事后}结算量；
第一阶段变量在全部情景间共享。式~\eqref{eq:q2-plan} 用\textbf{不等号}，即计划购电量只要求不低于
“预测口径净需求 \(+\) 充电量 \(-\) 放电量”，故允许超购；式~\eqref{eq:q2-gap} 的缺口由紧急购电补齐，
因 \(5p_t>0\)，最优解中 \(g_{t,\omega}\) 取其下界

\begin{equation}
  g_{t,\omega}=\max\Bigl\{0,\ \bigl(L_{\omega,t}-PV_{\omega,t}\bigr)\Delta t+c_t-d_t-E_t\Bigr\},
  \label{eq:q2-gapmax}
\end{equation}

即式~\eqref{eq:q2-gap} 的紧形式；其中的充放电量为\textbf{计划值} \(c_t,d_t\)（不含 \(\omega\)）。
主口径 \(K=60\) 下每个子问题有 \(576+144\times60=9216\) 个变量、144 条等式与
\(144\times61\) 条不等式，仍是纯 LP，用 HiGHS 精确求解。结果窗口内逐日顺序推进共 334 个子问题
（日初储电量取前一日末值，1 月作预热期，假设 A8），终端条件 \(S_{144}\ge S_0\)（假设 A6）。

\subsubsection{场景生成与削减}\label{sec:q2-scen}

光伏预报误差是本问唯一的不确定性来源。定义\textbf{逐日逐槽残差}

\begin{equation}
  r_{d,t}=\bigl(\widehat{PV}^{\text{fcst}}_{d,t}-PV^{\text{act}}_{d,t}\bigr)\Delta t,
  \qquad d\in\mathcal D,\ t\in\mathcal T\quad(\text{kWh}),
  \label{eq:q2-resid}
\end{equation}

样本为窗口内 \(|\mathcal D|\times144=334\times144\) 个。场景集按三层生成，其思路与聚类式典型样本
提取一致~\cite{reflixiang}：\textbf{① 按日总残差分层}——按 \(s_d=\sum_tr_{d,t}\) 升序把 334 天等频
划分为 \(K\) 层（主口径 \(K=60\)，每层 5--6 天），保留残差水平分布的\textbf{全部区间}，尾部不被
均值化抹平；\textbf{② 层内取 medoid}——每层取“到本层其他样本欧氏距离之和最小”的\emph{真实历史日}
残差向量（不中心化、不平滑），逐槽保留日内形状与时段间相关；\textbf{③ 权重取层样本占比}
\(W_\omega=n_\omega/334\)，\(\sum_\omega W_\omega=1\)。情景 \(\omega\) 下的实际净需求为
\(ND_{\omega,t}=ND^{\text{fcst}}_t+r_{\omega,t}\)，其中
\(ND^{\text{fcst}}_t=(L^{\text{fcst}}_t-\widehat{PV}^{\text{fcst}}_t)\Delta t\)。
场景集一次性生成后\textbf{存盘固定}，问题 3、4 与全部消融实验读同一份，以保证改进来自模型而非
场景噪声。三种层数下总费用极差 \(13.65\) 万元，仅为主口径的 \(0.96\%\)，\textbf{结论对层数稳健}；
紧急购电量对 \(K\) 更敏感（\(K=40\) 时为 \(K=120\) 的 1.68 倍）。\(K=1\)（残差恒为零、严格按预报
口径购电）即后文的\textbf{无安全边际}消融。

\subsubsection{报童型安全边际}\label{sec:q2-newsvendor}

暂不考虑储能的跨期耦合，把每槽视为独立报童问题~\cite{refzengming}：计划购电量 \(E_t\) 相当于订货量，
实际需求 \(D_t=(L_t-PV_t)\Delta t+c_t-d_t\) 相当于需求量；缺 1~kWh 要多付 \((5-1)p_t=4p_t$
（\textbf{缺货成本} \(c_u=4p_t\)，正常电价已计入计划口径），多买 1~kWh 的损失记为 \(c_o\)。
最优计划量是需求分布的\textbf{临界分位数}

\begin{equation}
  q^{*}_t=F_{D_t}^{-1}\!\left(\frac{c_u}{c_u+c_o}\right)
  =F_{D_t}^{-1}\!\left(\frac{4p_t}{4p_t+c_o}\right),
  \label{eq:q2-newsvendor}
\end{equation}

与之互补的是\textbf{最优缺货概率} \(1-q^*_t\)，两者不可混写。过剩成本 \(c_o\) 的两种口径给出
临界分位数的两条方向性参照：\textbf{口径一}（超购全废弃）\(c_o=p_t\)，得 \(q^*=4/5=0.800\)、
缺货概率 \(1/5=0.200\)；\textbf{口径二}（超购被储能吸纳并在后续低价时段替代购电）按最低可替代
购电价 \(p^{\min}=0.3713\)~元/kWh 计回收价值，\(c_o=p_t-p^{\min}\approx0.3949\)~元/kWh，取
\(p_t=0.7662\) 得 \(q^*\approx0.886\)、缺货概率 \(0.114\approx1/9\)。两者构成约
\([0.114,\,0.200]\) 的参考区间；它们来自单期近似，\textbf{不是严格上下界}（表~\ref{tab:newsvendor}）。

\begin{table}[htbp]
  \centering
  \caption{问题 2 主模型的逐日紧急购电量与报童两条方向性参照（对象：334 天逐日紧急购电量与 7 日滑动平均；
            两条水平线为缺货概率 \(0.200\) 与 \(0.114\) 对应的分位线；实测缺货频率 8.55\%、满足率 99.04\%）}
  \label{tab:newsvendor}
  \footnotesize
  \begin{tabularx}{\linewidth}{L{3.4cm} C{2.0cm} C{2.4cm} C{1.7cm} Y}
    \toprule
    过剩成本口径 & 临界分位数 & 最优缺货概率 & 样本量 & 与实测的关系 \\
    \midrule
    口径一：超购全废弃 \(c_o=p_t\) & 0.800 & 0.200（\(1/5\)） & --- & 参照上侧 \\
    口径二：超购可回收 \(c_o=p_t-p^{\min}\approx0.395\)~\unit{元/kWh} & 0.886 & 0.114（\(\approx1/9\)） & --- & 参照下侧 \\
    \midrule
    两阶段随机规划实测 & 覆盖率 0.9145 & \textbf{0.0855} & 48\,096 槽 & 低于两条参照（更保守） \\
    \bottomrule
  \end{tabularx}
  \\[1ex]
  {\footnotesize 表注：实测缺货频率 \(=\) 紧急购电量大于 \(10^{-6}\)~kWh 的槽数占 \(334\times144\) 槽的
  比例；覆盖率 \(=1-\) 缺货频率；电量满足率为 \(99.04\%\)（\calref{C-05}）。}
\end{table}

实测行为：逐槽缺货频率 \(8.55\%\)（覆盖率 \(91.45\%\)）\textbf{低于两条参照的下侧}，电量满足率
\(99.04\%\)，全窗口安全边际合计 \(576\,299.67\)~kWh（平均 11.98~kWh/槽）\evid{C-Q2-01}。
更保守的原因在结构而非参数：\textbf{储能同时承担对冲功能}，同一份计划可由电池放电在缺口时段补足，
单期报童模型无法刻画“计划购电量”与“电池吞吐”的替代关系。故式~\eqref{eq:q2-newsvendor} 与
表~\ref{tab:newsvendor} 只作\textbf{方向性近似与参照区间}，\textbf{不是}精确最优解，其适用边界有四条：
① \textbf{多时段耦合}——144 槽通过储电量递推耦合，本槽多买的电占用容量并影响其后所有槽，
不能拆成 144 个独立报童问题；② \textbf{储能缓冲}——过剩电量可被吸纳并在后续高价时段替代购电，
故 \(c_o<p_t\)，临界分位数高于口径一的 0.800；③ \textbf{紧急购电的时标在结算时刻}——缺口只能事后按
\(5p_t\) 补足，不能回溯改写计划；④ \textbf{残差分布未知}——以 334 天样本分位数替代总体分位数，
每槽仅 334 个观测，估计带误差，须报区间。

\textbf{安全边际的经济价值}由无安全边际消融量化：场景集退化为 \(K=1\) 后，全年计划购电费
\(12\,570\,533.91\) 元、紧急购电费 \(4\,828\,025.76\) 元，合计 \(17\,398\,559.68\) 元，
\textbf{比无储能基线 \(16\,407\,319.63\) 元还高 \(6.04\%\)}，紧急购电量升至 \(1\,156\,618.02\)~kWh
且 334/334 天全部触发。分解为：安全边际把计划购电费推高 \(1\,067\,018.99\) 元，把紧急购电费压低
\(4\,200\,786.52\) 元，\textbf{净收益 \(3\,133\,767.55\) 元}（相对不设边际降低 \(18.0\%\)）。
边际电量只占全窗口计划购电量的 \(2.59\%\)，却消掉 \(87.0\%\) 的紧急购电费。惩罚倍数的方向一致：
同一典型日与同一场景集上把倍数由 \(k=2\) 提高到 \(k=10\)，单日最优安全边际由 \(671\)~kWh 增至
\(3\,334\)~kWh，期望紧急购电量由 \(1\,246\)~kWh 降至 \(187\)~kWh（\(k=5\) 时为 \(428\)~kWh）。

\subsection{模型求解与结果分析}

\subsubsection{求解结果与复核}\label{sec:q2-results}
在 334 天窗口内求解式~\eqref{eq:q2}，总费用 \textbf{14\,264\,792.13 元}（1\,426.48 万元），其中
计划购电费 \(13\,637\,552.90\) 元、紧急购电费 \(627\,239.24\) 元，占比 \(95.60\%\) 与 \(4.40\%\)
（\calref{C-03}）；相对 \(C^{base}=16\,407\,319.63\) 元（变体 V2）节省 \(2\,142\,527.50\) 元，
\textbf{节省率 \(13.06\%\)}（\calref{C-04}）。全窗口计划购电量 \(22\,220\,333.20\)~kWh
（\calref{C-02}：模型内值，等于结果文件「全天购电量」汇总列之和；若按 144 个逐槽 4 位小数取整列求和，
则为 \(22\,220\,333.14\)~kWh，相对差 \(2.9\times10^{-9}\)——两者只差写入时的取整，非模型差异），
紧急购电量 \(172\,015.82\)~kWh、出现在 296/334 天（\calref{C-05}），
逐槽缺货频率 \(8.55\%\)、电量满足率 \(99.04\%\)、安全边际合计 \(576\,299.67\)~kWh。

\subsubsection{四个指定日期的结果表（题面表 1／表 2／表 3）}

按题面要求，问题 2 需对 \(2025\text{-}03\text{-}20\)、\(06\text{-}21\)、\(09\text{-}23\)、\(12\text{-}21\)
四个日期各给出表 1 格式（指定 \(6\) 个 \(10\)~min 时段的购电量、全天购电量与全天购电费）、
表 2 格式（\(6\) 个 \(4\)~h 块的充放电量与两端储电量）与表 3 格式（紧急购电时间段与购电量）。
四日均满足计划口径供需平衡（最大残差 \(1.82\times10^{-12}\)）与日周期条件。

\begin{table}[htbp]
  \centering
  \caption{问题 2 在四个指定日期的购电量（表 1 格式；单位：购电量 kWh、购电费 元；
           时段为 10~min，按题面模板的位置读法，\(10{:}00\text{-}10{:}10\) 对应附件标签 \(10{:}10\)；
           全天购电量与全天购电费为对应日期 144 个槽的合计）}
  \label{tab:q2-t1}
  \footnotesize
  \setlength{\tabcolsep}{5pt}
  \begin{tabular}{lcccc}
    \toprule
    时间段 & 2025-03-20 & 2025-06-21 & 2025-09-23 & 2025-12-21 \\
    \midrule
    \(10{:}00\)--\(10{:}10\) & 0.0000 & 0.0000 & 0.0000 & 0.0000 \\
    \(12{:}00\)--\(12{:}10\) & 804.3029 & 0.0000 & 536.2415 & 997.0335 \\
    \(14{:}00\)--\(14{:}10\) & 0.0000 & 0.0000 & 0.0000 & 0.0000 \\
    \(16{:}00\)--\(16{:}10\) & 817.3061 & 255.2254 & 902.7969 & 1\,001.2240 \\
    \(18{:}00\)--\(18{:}10\) & 697.9973 & 0.0000 & 764.3298 & 715.4144 \\
    \(20{:}00\)--\(20{:}10\) & 0.0000 & 0.0000 & 0.0000 & 0.0000 \\
    \midrule
    全天购电量（kWh） & 74\,178.1298 & 35\,701.5391 & 74\,744.7963 & 99\,692.6601 \\
    全天购电费（元） & 45\,229.90 & 20\,324.26 & 47\,063.39 & 64\,376.98 \\
    \bottomrule
  \end{tabular}
\end{table}

\begin{table}[htbp]
  \centering
  \caption{问题 2 在四个指定日期的储能充放电量与储电量（表 2 格式；单位：kWh；
           6 个 4~h 块各含 24 个 10~min 槽；\(0{:}00\) 与 \(24{:}00\) 储电量对应题面表 2 的“时刻/储电量”列，
           四个日期均为 \(6000.00\)~kWh）}
  \label{tab:q2-t2}
  \footnotesize
  \setlength{\tabcolsep}{4pt}
  \begin{tabular}{l*{4}{cc}}
    \toprule
    \multirow{2}{*}{时间段} & \multicolumn{2}{c}{2025-03-20} & \multicolumn{2}{c}{2025-06-21}
                            & \multicolumn{2}{c}{2025-09-23} & \multicolumn{2}{c}{2025-12-21} \\
    \cmidrule(lr){2-3}\cmidrule(lr){4-5}\cmidrule(lr){6-7}\cmidrule(lr){8-9}
     & 充电量 & 放电量 & 充电量 & 放电量 & 充电量 & 放电量 & 充电量 & 放电量 \\
    \midrule
    \(0{:}00\)--\(4{:}00\)   & 4\,500.0000 & 0.0000 & 0.0000 & 0.0000 & 4\,500.0000 & 0.0000 & 4\,500.0000 & 0.0000 \\
    \(4{:}00\)--\(8{:}00\)   & 833.3333 & 6\,224.1108 & 140.7851 & 4\,303.4721 & 833.3333 & 5\,547.8264 & 1\,666.6667 & 1\,508.3333 \\
    \(8{:}00\)--\(12{:}00\)  & 6\,500.0000 & 2\,415.8892 & 5\,207.8672 & 130.5638 & 5\,833.3333 & 3\,020.9227 & 5\,833.3333 & 7\,806.6667 \\
    \(12{:}00\)--\(16{:}00\) & 6\,905.5931 & 2\,218.5304 & 5\,458.7995 & 0.0000 & 6\,728.4081 & 1\,606.2615 & 8\,333.3333 & 2\,835.0000 \\
    \(16{:}00\)--\(20{:}00\) & 0.0000 & 5\,709.7835 & 0.0000 & 6\,154.1006 & 0.0000 & 5\,602.9091 & 0.0000 & 5\,797.4742 \\
    \(20{:}00\)--\(24{:}00\) & 5\,333.3333 & 2\,930.2165 & 5\,333.3333 & 2\,485.8994 & 5\,333.3333 & 3\,037.0909 & 5\,333.3333 & 2\,842.5258 \\
    \midrule
    \(0{:}00\) 储电量 & \multicolumn{2}{c}{6\,000.00} & \multicolumn{2}{c}{6\,000.00}
                      & \multicolumn{2}{c}{6\,000.00} & \multicolumn{2}{c}{6\,000.00} \\
    \(24{:}00\) 储电量 & \multicolumn{2}{c}{6\,000.00} & \multicolumn{2}{c}{6\,000.00}
                       & \multicolumn{2}{c}{6\,000.00} & \multicolumn{2}{c}{6\,000.00} \\
    \bottomrule
  \end{tabular}
\end{table}
\begin{table}[htbp]
  \centering
  \caption{问题 2 在四个指定日期的紧急购电时间段与购电量（题面表 3 格式；单位：kWh；
            时段按连续触发的 10~min 槽合并、左闭右开；区间值取 2 位、日合计由全精度求和后取 2 位，
            结果文件则逐区间存 4 位小数，故个别日合计与文件逐区间求和可差 \(1\times10^{-4}\)~kWh，
            审计项 AU10 容差 \(10^{-2}\)~kWh）}
  \label{tab:emg-q2}
  \footnotesize
  \setlength{\tabcolsep}{4pt}
  \begin{tabular}{C{1.85cm} L{9.30cm} C{1.85cm}}
    \toprule
    日期 & 紧急购电时间段（括号内购电量，kWh） & 合计（kWh） \\
    \midrule
    2025-03-20 & \(6{:}00\)--\(8{:}00\)（228.99）；\(8{:}10\)--\(8{:}30\)（49.16）；\(8{:}50\)--\(10{:}30\)（531.27）；\(10{:}40\)--\(11{:}40\)（216.46） & 1\,025.88 \\
    2025-06-21 & \(4{:}00\)--\(4{:}50\)（5.24）；\(14{:}00\)--\(14{:}30\)（181.84） & 187.08 \\
    2025-09-23 & \(6{:}50\)--\(7{:}40\)（132.59）；\(11{:}00\)--\(11{:}40\)（60.77） & 193.37 \\
    2025-12-21 & \(11{:}00\)--\(11{:}30\)（22.53） & 22.53 \\
    \bottomrule
  \end{tabular}
\end{table}


\subsection{检验与灵敏度分析}

\subsubsection{内部一致性与基线对比}
逐日复核：计划口径残差最大 \(3.41\times10^{-13}\)~kWh，储电量递推残差最大
\(1.09\times10^{-11}\)~kWh；同时充放、储电量越界、功率越界与负购电的时段数均为 0；末值低于初值的
天数 0；费用相对复算误差 \(\le10^{-6}\)；结果文件回读的全天购电量最大偏差 \(0.0011\)~kWh、
全天购电费最大相对偏差 \(3.40\times10^{-7}\)。

\textbf{基线对比。} 不做优化（无储能）时窗口购电费 \(C^{base}=16\,407\,319.63\) 元（变体 V2），
两阶段计划为 \(14\,264\,792.13\) 元，\textbf{下降 \(2\,142\,527.50\) 元（\(13.06\%\)）}（\calref{C-04}）。
更关键的是\textbf{安全边际这一步本身的边际价值}：取消安全边际（确定性等价、零残差情景）后总费用升至
\(17\,398\,559.68\) 元，比无储能基线还高 \(6.04\%\)，紧急购电量由 \(172\,015.82\) 升至
\(1\,156\,618.02\)~kWh（\(6.72\) 倍），故安全边际的边际价值为 \(3\,133\,767.55\) 元——
它是本问区别于问题 1 的核心模块（\S\ref{sec:q2-newsvendor}）。

\textbf{消融对照。} 把本问的模块逐一移除后的效果见表~\ref{tab:abl-q2}。\textbf{安全边际不是装饰}：整行移除后费用升至
\(17\,398\,559.68\) 元、比无储能基线高 \(6.04\%\)，紧急购电量由 \(172\,015.82\) 增至
\(1\,156\,618.02\)~kWh（6.72 倍），故“按情景分位数超购”是问题 2 的核心模块；与实测对照，
逐槽缺货频率 \(8.55\%\)、电量满足率 \(99.04\%\)，落在报童两条方向性参照 \(0.114\)--\(0.200\) 的\textbf{下侧}
（储能除承担安全边际外还起对冲作用），故报童结论只作方向性近似与参照区间 \evid{C-Q2-01}。

\begin{table}[htbp]
  \centering
  \caption{问题 2 的基线与消融对照（费用单位：元，电量单位：kWh；全部为 2025-02-01 至 12-31 共 334 天、
           附件 1 固定电价口径；“相对 \(C^{base}\)”以无储能基线 \(16\,407\,319.63\) 元
           （\calref{C-01} 变体 V2）为分母）}
  \label{tab:abl-q2}
  \scriptsize
  \setlength{\tabcolsep}{3pt}
  \begin{tabularx}{\linewidth}{C{0.72cm} L{3.15cm} C{2.05cm} C{1.55cm} C{1.65cm} Y}
    \toprule
    编号 & 配置（新增或移除的模块） & 总费用 & 相对 \(C^{base}\) & 紧急购电量 & 结论 \\
    \midrule
    B0 & 无储能：净负荷全额购电（闭式 \calref{C-01}）
      & 16\,407\,319.63 & 0 & --- & 基线 \\
    R2 & 典型日预报：计划依据改为附件 1 的日均形状
      & 23\,484\,646.63 & \(-43.14\%\) & 2\,525\,548.76 & 比无储能基线还贵 \(7\,077\,327.00\) 元，不可用 \\
    E0 & 无安全边际：确定性等价（\(K=1\)、零残差情景）
      & 17\,398\,559.68 & \(-6.04\%\) & 1\,156\,618.02 & 移除后反而高于无储能基线；
        \textbf{安全边际的边际价值 \(3\,133\,767.55\) 元} \\
    R1 & 完全信息：\(0{:}00\) 知当日真值
      & 13\,411\,560.90 & 18.26\% & 0.00 & 预测误差成本 \(853\,231.23\) 元（下界） \\
    M & 主模型：\(0{:}00\) 一次预报 \(+\) 两阶段随机规划（\(K=60\)）
      & \textbf{14\,264\,792.13} & 13.06\% & 172\,015.82 & 主方案（基准行） \\
    K120 & 场景数加倍（\(K=120\)，其余同主模型）
      & 14\,235\,163.95 & 13.24\% & 141\,707.94 & 费用仅降 \(29\,628.19\) 元（\(-0.21\%\)），结论稳健 \\
    RT & \textbf{实时再调度}（假设 A12 的消融）：\(E^{\text{plan}}\) 固定为已结算的计划量，
         用当日实际净需求重优化 \(c,d,S\)
      & 13\,839\,967.14 & 15.65\% & 62\,810.33 & 紧急购电降 \(63.49\%\)（费用降 \(67.73\%\)），
        但盈余处置量只降 \(1.12\%\)——见下段说明 \\
    \bottomrule
  \end{tabularx}
\end{table}

\textbf{充放电冻结假设 A12 的代价被量化}：把 \(E^{\text{plan}}\) 固定在
生产解（已按计划量结算、属沉没成本）、只让 \(c,d,S\) 按当日\emph{实际}净需求重优化后，
紧急购电量由 \(172\,015.82\) 降到 \(62\,810.33\)~kWh（\(-63.49\%\)）、紧急购电费由 \(627\,239.24\)
降到 \(202\,414.24\) 元（\(-67.73\%\)）、触发天数由 296 降到 188，总费用降 \(424\,825.00\) 元
（\(13\,839\,967.14\) 元，节省率 \(15.65\%\)）。\textbf{但盈余处置量几乎不动}
（\(3\,106\,390.84\to3\,071\,686.16\)~kWh，\(-1.12\%\)），原因是结构性的：由 A6 的终端条件
\(S_{143}\ge S_0\) 与往返效率 \(0.81\) 可得 \(\sum_td_t\le0.81\sum_tc_t\)，
\textbf{一天内电池不能净供能}，放电必须由同日充电补偿；因此 310 万 kWh 的处置量\textbf{主要来自
“按预报超购后实际用不上”}，而不是“电池不肯动”。三条边界：① 本行\emph{只}放开调度、未重新优化计划，
计划与调度联合重优化的收益未测（需展开多阶段模型）；② \(62\,810.33\)~kWh 是“计划固定下不可再减”
的最小值，故 \(63.49\%\) 是\textbf{该假设代价的下界}；③ 若进一步放开计划可行性行
\(E^{\text{plan}}_t-c_t+d_t\ge\widehat{ND}_t\)，残余紧急购电可再降到 \(14\,964.60\)~kWh，
但该行是本文模型的一部分，故保留。

\subsubsection{对照方案、机会约束与灵敏度分析}\label{sec:q2-compare}\label{sec:q2-chance}
在同一批实际轨迹上比较三种计划依据（表~\ref{tab:abl-q2} 给出完整口径）：

\begin{equation}
  \underbrace{13\,411\,560.90}_{\text{R1 完全信息（下界）}},\quad
  \underbrace{14\,264\,792.13}_{\text{R3 主方案}},\quad
  \underbrace{23\,484\,646.63}_{\text{R2 典型日预报（不可用）}}\quad(\text{元}),
  \label{eq:q2-compare}
\end{equation}

R1 的紧急购电量恒为 0，其费用是“预测误差成本”的下界；若放弃日决策解耦、把 334 天写成单一
LP，则得 \(12\,228\,166.56\) 元，是更松的下界。R2 比无储能基线还高 \(43.14\%\)（紧急购电量
\(2\,525\,548.76\)~kWh、286/334 天触发）——附件 1 是实际功率的\emph{全年逐时段均值}而非日前预报，
只用于说明“预报质量决定成本”。R3 与 R1 之差

\begin{equation}
  \Delta C_{\text{forecast}}=14\,264\,792.13-13\,411\,560.90=853\,231.23\ \text{元}
  \label{eq:q2-forecasterr}
\end{equation}

即\textbf{预测误差的成本}，也就是“5 倍紧急购电”机制对应的金额。

\textbf{机会约束对照。} 期望费用最小并不直接约束缺电概率，工程上另一种提法是要求逐槽

\begin{equation}
  \Pr\bigl(E_t-c_t+d_t\ \ge\ \widetilde{ND}_t\bigr)\ \ge\ 1-\beta ,
  \label{eq:q2-chance}
\end{equation}

其中 \(\beta\) 与式~\eqref{eq:q2-newsvendor} 的最优缺货概率同义，故本节把它作为\textbf{解析—数值
互证}~\cite{refzhangzt}。式~\eqref{eq:q2-chance} 含决策变量，直接求解需 0--1 变量，故用其保守确定性
等价：把安全边际设在情景净需求的 \(1-\beta\) 分位数上，即
\(E_t-c_t+d_t\ge q_{1-\beta}(\{\widetilde{ND}_{\omega,t}\})\)。该条件对式~\eqref{eq:q2-chance}
\textbf{充分}，但忽略储能的跨槽对冲，故\textbf{保守}。结果见表~\ref{tab:chance}。

\begin{table}[htbp]
  \centering
  \caption{机会约束对照（同一 334 天、同一实际轨迹、同一终端条件；安全边际取情景净需求的
           $1-\beta$ 分位数；缺电频率为 334$\times$144 个槽位中触发紧急购电的比例；
           费用单位元，CVaR$_{0.95}$ 为逐日费用的 95\% 条件风险价值，单位元/日）}
  \label{tab:chance}
  \small
  \setlength{\tabcolsep}{3pt}
  \begin{tabular}{crrrrrcr}
    \toprule
    $\beta$ & 总费用 & 计划购电费 & 紧急购电费 & 紧急购电量 & 缺电频率 & 逐日 CVaR$_{0.95}$ & 节省率 \\
    \midrule
    $0.01$ & 15\,010\,972.12 & 14\,957\,838.66 & 53\,133.46 & 12\,022.0 & 0.82\% & 70\,158.17 & 8.51\% \\
    $0.10$ & 14\,339\,249.60 & 13\,889\,543.70 & 449\,705.90 & 112\,236.5 & 6.21\% & 66\,086.57 & 12.60\% \\
    $0.20$ & 14\,417\,662.70 & 13\,572\,370.56 & 845\,292.14 & 222\,694.0 & 10.15\% & 65\,746.11 & 12.13\% \\
    $0.50$ & 15\,135\,214.64 & 13\,063\,582.93 & 2\,071\,631.71 & 563\,981.7 & 17.65\% & 68\,490.31 & 7.75\% \\
    \midrule
    \textbf{随机规划}（\ref{sec:q2-model}） & $\mathbf{14\,264\,792.13}$ & 13\,637\,552.90 & 627\,239.24 & 172\,015.82 & 8.55\% & $\mathbf{65\,663.71}$ & $\mathbf{13.06\%}$ \\
    \bottomrule
  \end{tabular}
\end{table}

表~\ref{tab:chance} 给出四条结论。① \textbf{总费用随 \(\beta\) 先降后升}：由 \(0.50\) 收紧到
\(0.10\) 下降 \(795\,965.04\) 元，再从 \(0.10\) 收紧到 \(0.01\) 反而上升 \(671\,722.52\) 元，
最优档 \(\beta=0.10\) 为 \(14\,339\,249.60\) 元。② \textbf{随机规划仍严格更优}：最优档仍贵
\(74\,457.47\) 元（\(+0.52\%\)），逐日 \(\mathrm{CVaR}_{0.95}\) 也更高（\(66\,086.57\) 对
\(65\,663.71\) 元/日），原因是分位数边际“逐槽各自为战”，而随机规划让储能做跨时段对冲。
③ \textbf{实测缺电频率低于标称 \(\beta\)}（\(\beta=0.10\) 时实测 \(6.21\%\)、\(0.20\) 时
\(10.15\%\)）。这里有\textbf{两个来源}不能混为一谈：其一是分位数边际“逐槽各自为战”、储能又额外承担
跨槽对冲，使实际缺电更少；其二是\textbf{情景分辨率的离散化误差}——60 个情景下经验分位数的概率分辨率
只有 \(1/60\approx1.67\%\)，\(\beta=0.10\) 的“\(90\%\) 分位”实际落在次序统计量的格点之间，标称
\(\beta\) 自身就带 \(\pm1.7\%\) 量级的取整误差，且情景分布与总体分布并不相同。本文\textbf{不把二者分离}，
故该条只作\textbf{方向性读数}，不外推为“储能对冲的保守度”。④ \textbf{报童界得到互证}：理论最优缺货
概率为 \(1/5=0.200\) 与约 \(0.114\) 两条方向性参照，而实测逐槽缺电频率 \(8.55\%\) \textbf{低于}两者。\evid{C-Q2-02}

\textbf{风险维度。} 期望费用最小化对极端情景并不敏感。在 334 天中每 6 天取 1 天的子样本（56 天）上
做 CVaR（\calref{C-11}）与 \(\lambda\) 扫描：\(\lambda\) 由 0 增至 1 时，期望费用由 \(41\,995.59\)
升至 \(42\,411.56\)~元/日（\(+0.99\%\)），而 \(\mathrm{CVaR}_{0.95}\) 由 \(50\,877.95\) 降至
\(49\,133.66\)~元/日（\(-3.43\%\)），日均紧急购电量由 \(480.2\) 降至 \(313.8\)~kWh。
\textbf{加权口径}：情景集是分层 medoid 集（权重 \(0.0030\)--\(0.0299\)），故上式 CVaR 按
\(w_\omega\) 加权、与目标函数中的期望项同权（\S\ref{sec:q4-frontier} 式~\eqref{eq:q4-cvar}）；
若误用等权平均，同一组解会报出 \(52\,176.36\to50\,510.46\)~元/日，量级相差约 \(2.5\%\)。
完整前沿与真实轨迹回测见 \ref{sec:q4-frontier}。

本问两个参数扰动对结论的影响见表~\ref{tab:sens-q2}。

\begin{table}[htbp]
  \centering
  \caption{问题 2 的参数灵敏度：扰动档位与“结论是否改变”的判断（费用单位 元，电量单位 kWh；
           单参数扫描除终端条件外均在同一组典型日参数下求解；完整档位实测值见
           \appref{sec:sens-app}{D.3 单参数灵敏度完整档位}）}
  \label{tab:sens-q2}
  \small
  \setlength{\tabcolsep}{4pt}
  \begin{tabularx}{\linewidth}{L{2.05cm} L{3.15cm} Y}
    \toprule
    参数 & 扰动档位 & 结论是否改变（含反转点） \\
    \midrule
    紧急购电倍数 \(\kappa\) & 2、3、5（基准）、8、10
      & 方向不变：\(\kappa\) 越大安全边际越大、缺口越小。反转点在 \(\kappa\to1\)——
        紧急购电与正常电价同价时缺货成本消失，“需要安全边际”的结论不再成立 \\
    场景数 \(K\) & 60（主口径）、120
      & 费用结论不变（仅降 \(0.21\%\)）；但紧急购电量降 \(17.62\%\)、触发天数由 296 降到 277——
        \textbf{分项指标对场景分辨率敏感，引用时必须带 \(K\)} \\
    \bottomrule
  \end{tabularx}
\end{table}
其中\textbf{场景数不是敏感维度}：\(K\) 由 60 加倍到 120 时费用只降 \(0.21\%\)，但紧急购电量降
\(17.62\%\)、触发天数由 296 降到 277——费用结论稳健，而该分项指标对场景分辨率敏感，引用时
必须带 \(K\)。

\textbf{假设 A9 的量化。} 题面只给光伏预报，未给负载预报，故 A9 把当日负载按已知处理，
此前的声明只是“会低估紧急购电量”。补上量级：只改\textbf{计划阶段}的负载信息，令
\(\widehat L_t=L_t+\varepsilon_t\)，\(\varepsilon\) 为按相对 MAE 标定的零均值误差
（\(\mathrm{MAE}_{\text{rel}}=\mathbb{E}|\varepsilon_t|/\overline L\)，\(\overline L=4626.03\)~kW），
模型、情景机制与结算口径全部不变，重跑 334 天（表~\ref{tab:loaderr}）。

\begin{table}[htbp]
  \centering
  \caption{负载预报误差对问题 2 的影响（只改计划期的负载信息；范围：2025-02-01 至 12-31 共 334 天；
           费用单位元、电量单位 kWh；样本：3 档相对 MAE；\(0\%\) 档与主结果逐位一致）}
  \label{tab:loaderr}
  \small
  \setlength{\tabcolsep}{5pt}
  \begin{tabular}{crrrrc}
    \toprule
    负载预报相对 MAE & 总费用 & 紧急购电量 & 出现天数 & 相对 \(C^{base}\) 节省率 & 相对 \(0\%\) 档费用 \\
    \midrule
    \(0\%\)（主口径） & \(14\,264\,792.13\) & \(172\,015.82\) & \(296\) & \(13.06\%\) & --- \\
    \(5\%\)           & \(15\,611\,126.88\) & \(293\,748.90\) & \(239\) & \(4.85\%\)  & \(+9.44\%\) \\
    \(10\%\)          & \(17\,472\,694.03\) & \(479\,838.20\) & \(202\) & \(-6.49\%\) & \(+22.49\%\) \\
    \bottomrule
  \end{tabular}
\end{table}

两条读数。① \textbf{量级}：相对 MAE 达 \(10\%\) 时问题 2 的总费用上升 \(22.49\%\)，且\textbf{相对无
储能基线的节省率转为负值}（\(-6.49\%\)）——在预报很差的微网里，两阶段计划加紧急购电的组合可能不如
不投运储能。这把 A9 的方向性声明变成可引用的区间：\(13.06\%\) 只适用于负载预报误差可忽略（或与光伏
残差同量级）的场景。② \textbf{一个反直觉但可解释的读数}：误差变大时紧急购电总量上升
（\(172\,015.82\to479\,838.20\)~kWh）而\textbf{触发天数反而下降}（\(296\to202\)）。原因不是计数口径，
而是模型的保守化：零均值误差使情景集变宽，计划层按更宽的分位数多购，“少而大”的缺口取代了“多而小”的
缺口。这也从反面说明主结果里 296 天触发、日均仅 \(581\)~kWh 的形态，是 \(0{:}00\) 光伏预报
\textbf{存在系统性时段偏差}（上午高估、下午低估，见 \S\ref{sec:data-err}）的结果，而非随机噪声的结果。
\evid{C-Q2-01}

\subsubsection{稳健性小结}
本问结论在扰动下的稳定域可明确划出。① \textbf{方向不变}：\(\kappa=2\)--\(10\)、场景数 \(K=60\to120\)
（费用仅降 \(0.21\%\)）、四种信息结构（R1/R2/R3）与机会约束四档 \(\beta\) 下，两阶段随机规划均优于
同口径的确定性计划。② \textbf{一处会反转的边界}：把计划期的\textbf{负载预报}误差放大到相对 MAE
\(10\%\) 时，总费用上升 \(22.49\%\)，\textbf{相对无储能基线的节省率转为 \(-6.49\%\)}——
即在预报很差的微网里“两阶段计划 \(+\) 紧急购电”可能不如不投运储能（表~\ref{tab:loaderr}）。
③ \textbf{引用纪律}：紧急购电量对场景分辨率敏感（\(K:60\to120\) 时降 \(17.62\%\)、触发天数
\(296\to277\)），该分项必须与 \(K\) 同时引用；费用类结论则不敏感。
